\writeSectionFrame{\presentationTitle}{Das Problem}
\begin{frame}{Unsere Fernbedienung}
	\begin{center}
 		\includegraphics[scale=0.4]{\BILDERPFAD/problem1.png}
 	\end{center}
\end{frame}

\note{
 	Wir möchten eine Fernbedienung für unser SmartHome implementieren. Im Moment hat die Fernbedienung nur eine einzige Taste. Und diese Taste soll ein Gerät in unserem Haus ansteuern.
}

\begin{frame}{Mögliches Klassendiagramm}
    \begin{tikzpicture}[scale=0.6, transform shape]
        \umlclass[type=interface, x=9, y=9]{Light}{}{
            +turnOn() \\
            +turnOff()
        }
    
        \umlclass[type=interface, x=6, y=9]{AlarmSystem}{}{
            +turnOff() \\
            +secure()
        }
        
        \umlclass[type=interface, x=3, y=9]{TV}{}{
            +turnOff() \\
            +turnOn()
        }
    
        \umlclass[type=interface, x=1, y=6]{OutdoorLights}{}{
            +turnOn() \\
            +turnOff()
        }
    
        \umlclass[type=interface, x=5, y=5]{StereoSystem}{}{
            +turnOn() \\
            +turnOff() \\
            +setCd() \\
            +setDvd() \\
            +setRadio()
        }
    
        \umlclass[type=interface, x=9, y=4]{GardenLights}{}{
            +setSundown() \\
            +setSunrise() \\
            +turnOn() \\
            +turnOff()
        }
    
        \umlclass[type=interface, x=3, y=1]{Ventilator}{}{
            +fast() \\
            +middle() \\
            +slow() \\
            +turnOff()
        }
    
        \umlclass[type=interface, x=6, y=1]{Sprinkler}{}{
            +waterOn() \\
            +waterOff()
        }
    \end{tikzpicture}
\end{frame}

\note{
 	Dies sind die Klassen der Gerätehersteller unseres smarten Homes. Und wie zu erwarten haben sich die Hersteller auch nicht bemüht, einen einheitlichen Satz an Schnittstellen zur Verfügung zu
 	stellen. Unsere Fernsteuerung muss also in der Lage sein, die Geräte aller Hersteller anzusteuern unabhängig von den Unterschieden in den APIs. Und wenn neue Geräte hinzukommen, müssen diese auch eingebaut werden können. Dabei soll die Fernbedienung so wenig wie möglich über die einzelnen Geräte wissen.
}